\section*{Bab \ref{ch:g6:exploring-our-environment} --- Menjelajahi Lingkungan Kita}

\begin{solution}{exo:g6:exploring-our-environment:1}
\omterm{def:g6:exploring-our-environment:components}{Makhluk hidup} (\omterm{prop:g4:plants-without-seeds:damp}{lumut}, kutu kayu); \omterm{def:g6:exploring-our-environment:components}{materi mineral} (batu temboknya,
genangannya); jejak dan sisa \omterm{def:g6:exploring-our-environment:components}{makhluk hidup} (\omterm{def:g1:plants-around-us:parts}{daun} kering, \omterm{def:g1:animals-around-us:covering}{cangkang}
siputnya).
\end{solution}

\begin{solution}{exo:g6:exploring-our-environment:2}
\omterm{def:g6:exploring-our-environment:conditions}{Suhu}, dengan termometer (dalam derajat Celsius); \omterm{def:g6:exploring-our-environment:conditions}{kelembapan} udara dan
tanahnya; \omterm{def:g6:exploring-our-environment:conditions}{cahaya}, dari matahari penuh sampai naungan yang pekat.
\end{solution}

\begin{solution}{exo:g6:exploring-our-environment:3}
Hidup: \omterm{prop:g4:plants-without-seeds:damp}{tumbuhan pakunya}, kelabangnya. Mineral: genangannya, batu
temboknya. Jejak: \omterm{def:g1:animals-around-us:covering}{cangkang} siputnya, tinja \omterm{prop:g3:sorting-living-things:groups}{burungnya}.
\end{solution}

\begin{solution}{exo:g6:exploring-our-environment:4}
Enam belas derajat (\qty{38}{\degreeCelsius} melawan
\qty{22}{\degreeCelsius}) dan matahari penuh melawan naungan yang
pekat (juga kering melawan lembap). Rumput liar yang tahan kekeringan
dan silau cocok dengan kaki tembok selatan; \omterm{prop:g4:plants-without-seeds:damp}{lumut}, \omterm{prop:g4:plants-without-seeds:damp}{tumbuhan paku} dan
kutu kayu cocok dengan sudut yang sejuk dan lembap.
\end{solution}

\begin{solution}{exo:g6:exploring-our-environment:5}
Penutupnya membiarkan air lolos, sehingga di bawah matahari yang
kering tubuhnya kehilangan airnya secepat itu sampai mati. Itulah
sebabnya kutu kayu ditemukan di \omterm{def:g6:exploring-our-environment:microhabitat}{mikrohabitat} yang lembap dan gelap ---
di bawah batu, di bawah kayu mati.
\end{solution}

\begin{solution}{exo:g6:exploring-our-environment:6}
Misalnya: sebutir \omterm{prop:g3:flowers-fruits-seeds:transform}{buah} hazel yang tergigit --- tikus hutannya; sebuah
jaring laba-laba --- laba-labanya; tapak di lumpur --- kuntulnya
(tinja, \omterm{def:g1:animals-around-us:covering}{bulu burung} dan \omterm{def:g1:animals-around-us:covering}{cangkang} berlaku dengan cara yang sama).
\end{solution}

\begin{solution}{exo:g6:exploring-our-environment:7}
Sebuah kotak tertutup, separuhnya kertas lembap dan separuhnya kertas
kering, kutu kayunya dilepaskan di sepanjang garis tengahnya; beberapa
menit kemudian nyaris semuanya duduk di sisi yang lembap (demikian
pula naungan melawan \omterm{def:g6:exploring-our-environment:conditions}{cahaya}). Ini menunjukkan bahwa tiap \omterm{def:g1:animals-around-us:animal}{hewan}
menetap di tempat \omterm{def:g6:exploring-our-environment:conditions}{kondisinya} sesuai dengan tubuhnya --- peta
kelompoknya menggambar dirinya sendiri tanpa penggiringan.
\end{solution}

\begin{solution}{exo:g6:exploring-our-environment:8}
Karena hanya pengukuran yang dibuat dengan cara yang sama yang dapat
dibandingkan: sebuah perbedaan harus datang dari lokasinya, bukan dari
pengukurannya. Inilah versi lapangan dari uji yang adil --- satu hal
yang berubah (lokasinya), segala hal lain dijaga tetap setara.
\end{solution}

\begin{solution}{exo:g6:exploring-our-environment:9}
Bagian atas batunya: tersinari, kering, hangat --- \omterm{prop:g4:plants-without-seeds:damp}{lumut} kerak (dengan
lalat yang berjemur sebagai tamunya). Sisi bawah batunya: gelap,
lembap, sejuk --- kutu kayu (atau kelabangnya). Berjarak lima
sentimeter, dua perangkat penghuni yang berpasangan.
\end{solution}

\begin{solution}{exo:g6:exploring-our-environment:10}
Langkah 1: tandailah satu meter persegi yang pasti. Langkah 2:
ukurlah --- bacaan termometer dituliskan, \omterm{def:g6:exploring-our-environment:conditions}{kelembapan} tanahnya diraba
lalu dicatat, kelas \omterm{def:g6:exploring-our-environment:conditions}{cahayanya} dicatat. Langkah 3: carilah dengan
benar, batunya diangkat lalu dikembalikan, ketiga macam komponennya
didaftarkan. Langkah 4: catatlah di tempat --- hitungan, bukan
``beberapa''. Langkah 5: dengan satu lokasi saja tidak ada
perbandingan yang mungkin; surveilah lokasi kedua dengan cara yang
sama.
\end{solution}

\begin{solution}{exo:g6:exploring-our-environment:11}
Penjelasannya: tanaman ivi tumbuh subur dalam rentang \omterm{def:g6:exploring-our-environment:conditions}{kondisi} tembok
utara (naungan, \omterm{def:g6:exploring-our-environment:conditions}{kelembapan}) dan tidak dalam rentang tembok selatan
(silau, kekeringan). Ujinya: ukurlah kedua kaki temboknya --- \omterm{def:g6:exploring-our-environment:conditions}{suhu},
\omterm{def:g6:exploring-our-environment:conditions}{kelembapan}, \omterm{def:g6:exploring-our-environment:conditions}{cahaya} --- lalu periksalah kesukaan ivi yang sudah
diketahui; atau temukan ivi di tembok lain lalu catatlah arah hadap
apa yang dimiliki bersama oleh tembok-tembok itu.
\end{solution}

\begin{solution}{exo:g6:exploring-our-environment:12}
Singkirkanlah perbedaan dindingnya: jalankan pilihannya di dalam kotak
yang bundar (dindingnya sama di mana-mana), dengan separuh yang lembap
dan yang kering ditukar di antara percobaannya --- lembap di timur
pada percobaan 1, lembap di barat pada percobaan 2. Kalau kutu kayunya
mengikuti \omterm{def:g6:exploring-our-environment:conditions}{kelembapannya} ke mana pun ia diletakkan, dindingnya
tersingkir; hanya \omterm{def:g6:exploring-our-environment:conditions}{kondisi} yang diuji yang berubah di antara
percobaannya, segala hal lain dijaga tetap setara.
\end{solution}

\begin{solution}{pb:g6:exploring-our-environment:1}
\textbf{1.} Pagar tanamannya yang tinggi: ia menaungi parit N (dan
melindunginya), sedangkan parit S menghadap matahari sepanjang hari.

\textbf{2.} \omterm{def:g6:exploring-our-environment:conditions}{Suhu}: \qty{19}{\degreeCelsius} melawan
\qty{31}{\degreeCelsius}. \omterm{def:g6:exploring-our-environment:conditions}{Kelembapan}: tanah parit N lembap (melekat),
tanah parit S kering (debu). \omterm{def:g6:exploring-our-environment:conditions}{Cahaya} melengkapi ketiganya: naungan
melawan matahari penuh.

\textbf{3.} \omterm{def:g6:exploring-our-environment:conditions}{Kondisi} berubah sepanjang hari: menjelang pukul enam sore
setiap lokasi menjadi lebih sejuk dan \omterm{def:g6:exploring-our-environment:conditions}{cahayanya} lebih rendah, sehingga
perbedaan yang terukur akan mencampur lokasi dengan jam. Jam yang sama
menjaga perbandingannya tetap adil.

\textbf{4.} Supaya usaha pencariannya setara: daerah yang lebih luas
atau waktu yang lebih lama di satu lokasi akan menemukan lebih banyak
penghuni di situ --- sebuah uji yang tidak adil tentang parit mana
yang lebih kaya.

\textbf{5.} Kedua daftarnya hanya menyebutkan \omterm{def:g6:exploring-our-environment:components}{makhluk hidup}. Yang
hilang: komponen mineralnya (tanah, batu, air) dan jejaknya --- yang
juga harus dicatat pada bagian pertama sebuah pendataan yang benar.

\textbf{6.} Kutu kayu: parit N --- mereka membutuhkan \omterm{def:g6:exploring-our-environment:conditions}{kelembapan} dan
naungan atau mereka mengering. Kodoknya: parit N --- \omterm{def:g1:the-five-senses:sense}{kulitnya} lembap,
sehingga membutuhkan tempat berlindung yang lembap. Kadalnya: parit S
--- ia menghangatkan tubuhnya di bawah matahari, dan berburu di tempat
ia dapat berjemur.

\textbf{7.} Entah siput memang tinggal di situ dengan bersembunyi
(barangkali hanya giat pada jam malam yang lembap), atau \omterm{def:g1:animals-around-us:covering}{cangkangnya}
adalah sisa dari tempat lain atau musim lain --- sebuah jejak
membuktikan kehadiran \emph{pada suatu waktu}, bukan tempat tinggal
sekarang.

\textbf{8.} Gelap, lembap, sejuk. Aturan pada
\cref{prop:g6:exploring-our-environment:decide}: tiap tempat dihuni
oleh \omterm{def:g5:biodiversity:species}{spesies} yang rentangnya cocok dengan \omterm{def:g6:exploring-our-environment:conditions}{kondisinya}.

\textbf{9.} N punya \omterm{prop:g4:plants-without-seeds:damp}{lumut} dan kodok karena naungannya menjaganya tetap
sejuk dan lembap, yang memang dituntut tubuh mereka; S punya kadal dan
rumput liar berdaun lilin karena matahari penuhnya membuatnya panas
dan kering, yang cocok bagi seorang pemburu yang berjemur dan sebuah
\omterm{def:g1:plants-around-us:plant}{tumbuhan} yang tahan kekeringan.

\textbf{10.} Tanpa pagar tanamannya, parit N mendapat matahari penuh:
lebih panas, lebih kering --- \omterm{def:g6:exploring-our-environment:conditions}{kondisi} yang bergeser menuju parit S
hari ini. Harapkanlah para pencinta \omterm{def:g6:exploring-our-environment:conditions}{kelembapan} (\omterm{prop:g4:plants-without-seeds:damp}{lumut}, \omterm{prop:g4:plants-without-seeds:damp}{tumbuhan paku},
kutu kayu, kodoknya) menyusut dan para pencinta matahari (rerumputan
kering, belalang, barangkali kadalnya) pindah masuk.

\textbf{11.} Karena survei ulangnya harus berbeda dari yang pertama
dalam satu hal saja --- ketiadaan pagar tanamannya. Lokasi, alat atau
musim yang baru akan menambahkan perbedaannya sendiri, dan
perubahannya tidak lagi dapat ditimpakan pada pagar tanamannya: aturan
uji yang adil pada skala sebuah bentang alam.

\textbf{12.} Tiap kutu kayu mengembara dan berbelok sampai tubuhnya
merasa nyaman, sehingga kelompoknya berakhir berkumpul di tempat
\omterm{def:g6:exploring-our-environment:conditions}{kelembapannya} berada --- kedudukan \omterm{def:g1:animals-around-us:animal}{hewannya} mencatat \omterm{def:g6:exploring-our-environment:conditions}{kondisinya},
sepasti bacaan termometer. Tidak seorang pun menggiring mereka;
kenyamanan tiap individu sendirilah yang mengerjakan penggambaran
petanya, dan itulah persis sebabnya penghuni dapat dibaca sebagai
pengukuran.
\end{solution}
